% ============================================================
%  date/main.tex -- single-file working copy for the LENGTH CUT
%
%  Created 2026-09-16 from the v7 tree one level up
%  ("00 - DATE CAM Monitoring.tex" + section files 01-08), with every
%  \input inlined verbatim in order. Section boundaries are marked
%  with "% ===== NN - <file> =====" separators so edits can be traced
%  back to the source file if needed. The v7 tree is the archive; the
%  version history and change-log lives in its 00 file.
%
%  Target: DATE 2027 regular paper -- 6 pages + 1 page for references
%  only (CFP wording: "must not exceed 6 pages in length, with one
%  extra page allowed only for bibliographic references"). Final
%  paper deadline: Sunday 20 Sep 2026 AoE. Baseline at creation:
%  16 pages (~14.7 body + ~1.3 references).
%
%  STATUS: R1 structural cuts applied Sep 16 -> 14 pages clean. See the
%  length-cut log at the end of this file (With Comments build) and
%  cut_plan.md. backup/ holds the pre-round snapshots.
%
%  Supporting files copied alongside: IEEEtran.cls, IEEEtran.bst,
%  RobStd.sty, comments.tex (flag system), refs.bib, figures/.
%  Compile: bash compile.sh   (main.pdf = clean; main-comments.pdf)
%  or plain: pdflatex main && bibtex main && pdflatex main && pdflatex main
% ============================================================
\documentclass[conference]{IEEEtran}
\IEEEoverridecommandlockouts

\usepackage{cite}
\usepackage{amsmath,amssymb,amsfonts}
\usepackage{graphicx}
\usepackage{textcomp}
\usepackage[dvipsnames]{xcolor}
\usepackage{RobStd}
\usepackage{subcaption}
\usepackage{siunitx}
% --- siunitx v2/v3 compatibility shim ------------------------
\ifdefined\qty\else
  \NewDocumentCommand{\qty}{O{}mm}{\SI[#1]{#2}{#3}}
\fi
\ifdefined\unit\else
  \NewDocumentCommand{\unit}{O{}m}{\si[#1]{#2}}
\fi
% -------------------------------------------------------------
\usepackage{placeins}
\usepackage{multirow}
\usepackage{makecell}
\usepackage{booktabs}
\usepackage{url}
\renewcommand\arraystretch{1.15}
\setlength{\tabcolsep}{4pt}
\graphicspath{{figures/}}

\def\BibTeX{{\rm B\kern-.05em{\sc i\kern-.025em b}\kern-.08em
    T\kern-.1667em\lower.7ex\hbox{E}\kern-.125emX}}

% --- shorthand macros ----------------------------------------
\newcommand{\cam}{{CAM}\xspace}
\newcommand{\tcam}{{TCAM}\xspace}
\newcommand{\mcam}{{MCAM}\xspace}
\newcommand{\acam}{{ACAM}\xspace}
\newcommand{\imc}{{CiM}\xspace}
\newcommand{\fefet}{{FeFET}\xspace}
\newcommand{\fefets}{{FeFETs}\xspace}
\newcommand{\auroc}{{AUROC}\xspace}
\newcommand{\signrp}{{sign-RP}\xspace}
\newcommand{\tocite}[1]{[{\color{red}citation-#1}]}

% ============================================================
%  comments.tex -- annotation / review-flag system
%
%  STREAMLINED FLAG NAMES (Aug 2026, per MN: "this is just less typing"):
%      \dt[note]{content}   was \dothis   / \dothisflag
%      \eo[note]{content}   was \expandon / \expandonflag
%  The old names are kept as aliases below, so every existing .tex
%  still compiles unchanged. Semantics are identical.
%
%  All flags obey the showcomments boolean; compile.sh flips it.
% ============================================================
\usepackage{ifthen}
\usepackage{amssymb}
\usepackage{xcolor}
\usepackage{etoolbox}
\usepackage{environ}
\usepackage[most]{tcolorbox}

\newboolean{showcomments}
\setboolean{showcomments}{false}    % compile.sh targets this exact line -- keep it here

% DEFAULT IS FALSE (no comments) so Overleaf produces the clean copy.
% Run `bash compile.sh` to build both versions.

\definecolor{flagorange}{RGB}{200, 95, 0}
\definecolor{flagblue}{RGB}{0, 70, 190}
\definecolor{aiviolet}{RGB}{120, 40, 180}
\definecolor{dogreen}{RGB}{0, 130, 60}

% ---- the two action flags -----------------------------------
%  Uniform shape:  \flag[what to do]{content}
%    \dt  orange -- an action to take
%    \eo  blue   -- say more about this
%  Content survives into BOTH builds; the bracketed note appears only
%  in the With Comments build. With NO optional argument the braced
%  text IS the note and disappears from the clean build entirely:
%      \dt{fill in the submission date}          note only
%      \dt[Trim this]{The library is short ...}  keep text, flag it
%      \eo[Add numbers]{}                        note only
\newcommand{\flagmark}[3]{%
  \ifthenelse{\boolean{showcomments}}{{\color{#1}\bfseries[#2:\ #3]}\ }{}%
}
\newcommand{\dt}[2][]{%
  \ifstrempty{#1}%
    {\ifthenelse{\boolean{showcomments}}{{\color{flagorange}\bfseries[DO THIS:\ #2]}}{}}%
    {\flagmark{flagorange}{DO THIS}{#1}#2}%
}
\newcommand{\eo}[2][]{%
  \ifstrempty{#1}%
    {\ifthenelse{\boolean{showcomments}}{{\color{flagblue}\bfseries[EXPAND:\ #2]}}{}}%
    {\flagmark{flagblue}{EXPAND}{#1}#2}%
}
% ---- aliases (do not "clean these up") ----------------------
\let\dothis\dt
\let\dothisflag\dt
\let\expandon\eo
\let\expandonflag\eo

% ---- reviewer remarks --------------------------------------
% Inline remark:  \inote{short note}  /  \inote[aiviolet]{note}
\newcommand{\inote}[2][red]{%
  \ifthenelse{\boolean{showcomments}}{\textcolor{#1}{[#2]}}{}%
}
% Block remark. Escape underscores as \_ inside.
%   \begin{pnote} ... \end{pnote}    or    \begin{pnote}[aiviolet] ... \end{pnote}
\NewEnviron{pnote}[1][red]{%
  \ifthenelse{\boolean{showcomments}}{\begingroup\color{#1}\BODY\endgroup}{}%
}
% Review-build-only content (whole paragraphs, tables, or a page).
\newcommand{\reviewonly}[1]{\ifthenelse{\boolean{showcomments}}{#1}{}}

% ---- structural boxes (always visible) ---------------------
\newtcolorbox{sectionbox}{%
  colback=black, coltext=white,
  width=\linewidth, boxrule=0pt,
  left=3pt, right=3pt, top=2pt, bottom=2pt,
  fontupper=\sffamily\bfseries\normalsize,
  sharp corners,
  before=\par\vspace{4pt}, after=\par\vspace{2pt}}
\newtcolorbox{subsectionbox}{%
  colback=lightgray, coltext=black,
  width=\linewidth, boxrule=0pt,
  left=3pt, right=3pt, top=1pt, bottom=1pt,
  fontupper=\sffamily\bfseries\small,
  sharp corners,
  before=\par\vspace{3pt}, after=\par\vspace{1pt}}

% ---- per-author inline notes -------------------------------
\makeatletter
\newcommand{\mynote}[3]{%
  \ifthenelse{\boolean{showcomments}}{%
   \fbox{\bfseries\sffamily\scriptsize#1}%
   {\small$\blacktriangleright$\textsf{\emph{\color{#3}{#2}}}$\blacktriangleleft$}}%
  {\@bsphack\@esphack}%
}
\makeatother
\definecolor{PersianGreen}{rgb}{0, 0.65, 0.58}
\newcommand{\mn}[1]{\mynote{MN}{#1}{violet}}
\newcommand{\ib}[1]{\mynote{IB}{#1}{PersianGreen}}
\newcommand{\sj}[1]{\mynote{SJ}{#1}{blue}}
\newcommand{\xh}[1]{\mynote{XH}{#1}{orange}}

% \setboolean{showcomments}{true}   % or run: bash compile.sh

\begin{document}

% ============================================================
%  ANNOTATION FLAGS -- quick reference
%
%  All flags compile away when showcomments is false (the default).
%  The two action flags take the SAME shape, \flag[what to do]{content}.
%  Content survives into BOTH builds; the bracketed note appears only in
%  the With Comments build. With NO optional argument the braced text IS
%  the note and disappears from the clean build entirely.
%
%  ORANGE -- an action to take
%    \dt[note]{content}          (was \dothis / \dothisflag; still work)
%    \dt{fill in the deadline}   note only, vanishes from the clean build
%
%  BLUE -- say more about this
%    \eo[note]{content}          (was \expandon / \expandonflag; still work)
%
%  RED / VIOLET -- reviewer remarks, review build only
%    \inote{short note}   \inote[aiviolet]{AI-generated suggestion}
%    \begin{pnote} ... \end{pnote}     block note; escape _ as \_
%    \reviewonly{...}                  content only in the review build
%
%  PER-AUTHOR inline notes: \mn{...} \ib{...} \sj{...} \xh{...}
%
%  A ']' typed inside a flag's optional argument silently truncates it
%  and leaks the rest into the CLEAN build. Brace them: {[1]}, {[5]}.
% ============================================================

% ============================================================
%  WRITING CONVENTIONS FOR THIS DRAFT (MN's voice; fix-ai-writing skill)
%
%  1.  Plain, first person plural, hedged. "may", "could", "appears to",
%      "remains to be established". No absolutes that are not absolute.
%  2.  En dashes (--), never em dashes (---).
%  3.  Numbered inline lists: {\bf (1)} ... {\bf (2)} ..., not
%      "First, ... Second, ...".
%  4.  Bold run-in lead-ins end with a COLON, not a period, and the next
%      word keeps its capital:  \textbf{What the grid measures:} Whether ...
%  5.  Never the word "honest" as framing; never announce your own
%      fairness ("to be fair", "stated fairly", "candidly").
%  6.  No self-referential framing: "the whole argument", "the key
%      insight", "which is the entire problem", "the upshot", "X turns
%      on Y". State the substance and stop.
%  7.  No verdict tags on your own description ("and it works").
%  8.  Do not narrate design choices ("deliberately", "chosen to be").
%  9.  A section discusses; it is not a proof. "Sec. IV covers ...",
%      not "Sec. IV is the proof that ...".
%  10. No "binds" / "bites" for constraints; no borrowed-metaphor verbs
%      ("is the currency", "prices the choice", "lands on", "rests on").
%      Use "maps to", "depends on", "the limit here is".
%  11. American spelling, always. Define every acronym at first use.
%  12. Scope caveats go in footnotes, not into the argumentative flow.
%  13. Label invented or illustrative examples as such, in the text.
%  14. Figures follow the house-graphs conventions: Helvetica, Set1
%      colors, transparent background, provenance header in the
%      generating script (see gen_paper_figs.py).
% ============================================================

\bstctlcite{IEEEexample:BSTcontrol}

% DATE is double-blind: authors and funding are anonymized for
% submission. Funding restored at camera-ready (SRC/DARPA JUMP 2.0;
% NSF IRES; etc.).
% Title per MN (08-25): option (2) from the v02 markup -- names the
% design-space contribution first for a DATE audience.
\title{\cam-Resident Behavioral Monitoring for LLMs: Encoding Choices,
Cross-Domain Fidelity, and Library-Size Scaling}

\author{\IEEEauthorblockN{Anonymized for review}
\IEEEauthorblockA{DATE 2027 submission}}

\maketitle

\begin{abstract}
Frontier language model deployments increasingly treat a model's own
activations as a safety signal: a production jailbreak defense from
Anthropic scores a single linear probe over residual-stream
activations on every generated token, at roughly $10^{-3}$~percent
of forward-pass compute. What is monitored is set by $K$, the number
of stored behavioral signatures each activation is compared against,
and target $K$'s are increasing rapidly: from $K=1$ in that defense,
to $10^{2}$--$10^{3}$ once probes are fitted per behavior, domain,
and layer, to $1.3\!\times\!10^{5}$ per layer for a whole-vocabulary
readout. On a GPU the signature bank streams from high-bandwidth
memory once per monitored token, so at high enough $K$ or token rates
the monitor's own memory traffic can exceed the bandwidth of the
accelerator it runs on. We propose supporting this search with
content-addressable memories (\cam{}s), which evaluate distance
functions inside the memory array, trading arithmetic precision and
distance-function fidelity for orders of magnitude less data
movement. We quantify this trade on two workloads: {\bf (1)} a
ten-domain cross-domain deception grid over Llama-3.3-70B
activations, where 384-bit codes hold in-domain \auroc within 0.01
of the full-precision cosine ceiling at more than $300\times$
compression per signature; and {\bf (2)} whole-vocabulary readout,
where a 3-bit Euclidean code at $\sim$$10^{3}$ bits per row
reproduces 0.96 of the exact readout's top-1 selections. The
whole-vocabulary bank then occupies 6--\qty{33}{\mega\byte} per
layer, and its estimated in-array search energy is five to six
orders of magnitude below streaming the equivalent bank from HBM.
\end{abstract}

\begin{IEEEkeywords}
Content-addressable memory; compute-in-memory; large language models;
activation probes; AI safety monitoring; nearest-neighbor search;
ferroelectric FETs; Hamming distance
\end{IEEEkeywords}

%%%%%%%%%%%%%%%%%%%%%%%%%%%%%%%%%%%%%%%%%%%%%%%%%%%%%%%%%%%%%%%%%%%%%%
% ===== 01 - Introduction.tex =====
%%%%%%%%%%%%%%%%%%%%%%%%%%%%%%%%%%%%%%%%%%%%%%%%%%%%%%%%%%%%%%%%%%%%%%
\section{Introduction}
\label{sec:Introduction}
Large language model (LLM) deployments are wrapped in layered
safeguards, and until recently those layers read
\emph{text}~\cite{metr2026policies}. That has started to change: the
always-on stage of one production jailbreak defense now reads the
model's activations directly~\cite{ccpp2026,anthropic2026ccblog}, as
do deployed residual-stream probes for
Gemini~\cite{gemini2026probes}. Each of these is a single stored
vector scored once per generated token against one threshold, fitted
to one behavior. Agentic deployments appear to need more: in a recent
cyber evaluation campaign, evaluators identified 19 unsanctioned agent
actions, among them deception and the supply-chain compromise of a
real open-source project, from transcripts after the runs had
completed~\cite{aisi2026cyber,aisi2026incident}. One stored vector may
not cover such a setting, since probes transfer poorly across
behaviors: a probe fitted on factual lying scores near chance on
sycophancy~\cite{ying2026spectrum}, and probes trained on outright
lying miss deception carried by technically true
statements~\cite{berger2026liedetector}.

Coverage is therefore set by $K$, the number of stored signatures
each activation is compared against, and the published record
already spans five orders of magnitude: $K=1$ in the deployed
monitors, $K\approx10^{2}$--$10^{3}$ for probes fitted per behavior,
domain, and layer~\cite{ying2026spectrum,aisi2026incident}, and
$K\approx1.3\!\times\!10^{5}$ for a whole-vocabulary readout, in
which every token in the vocabulary has an internal pattern meaning
that word is currently active, whether or not it is ever
said~\cite{anthropic2026workspace}. On a GPU every stored signature
is read from high-bandwidth memory (HBM) once per monitored token, so
the monitor's traffic grows linearly in $K$ while its answer stays a
few bits: at $10^{3}$~tokens/s monitoring alone consumes an H100 at
$K\!\approx\!204$K and a Jetson-class part at $K\!\approx\!17$K
(Fig.~\ref{fig:ladder}, Sec.~\ref{subsec:CostModel}). At small $K$
the reads hide under inference latency, but the bytes and the energy
still move on every token of agent runs that reach $10^{8}$
tokens~\cite{aisi2026cyber}.

\begin{figure}[!t]
\centering
\includegraphics[width=3.4in]{fig_ladder_v3.pdf}
\caption{Library size $K$ at which monitoring traffic alone
($K\,d\,b\,L_{\mathrm{mon}}\,R$ bytes/s, Eq.~(\ref{eq:gputraffic}),
here at $d=8{,}192$, fp16, one monitored layer) reaches the memory
bandwidth of an H100, an A100, and a Jetson Thor, against token
rate.}
\label{fig:ladder}
\end{figure}

Three observations shape the design. {\bf (1)} The operation is
library search -- one query against $K$ stored rows, returning a
winner, a count, or the rows under a threshold -- which is what a
content-addressable memory (\cam) computes in place: the bank never
moves, and per-token traffic is flat in $K$. {\bf (2)} The array's
metric is not the detector's. \cam cells buildable in any CMOS
process report only whether a position disagrees, so a row measures
Hamming distance over a few hundred bits, while the detector is
defined over cosine or Euclidean distance on 8{,}192 floating-point
values; we therefore encode activations into a ternary code and a
thermometer code, and ask what a multi-bit emerging-device
cell~\cite{kazemi2021fefetcam} buys over a plain Hamming array at more
than $300\times$ compression. {\bf (3)} A trained stage in front of
the array can learn the classification itself, so fidelity alone does
not show that the detector is using its memory; the projections here
are frozen, and every configuration is audited with its stored rows
replaced by random vectors.

We evaluate these encodings at both ends of the $K$ range. On a
ten-domain deception grid over Llama-3.3-70B activations, 384-bit
codes hold in-domain \auroc (area under the receiver operating
characteristic curve) within 0.01 of the full-precision ceiling and
cost 0.065--0.126 across domains, a gap that no wider code,
multi-bit cell, or pooled bank closes. On whole-vocabulary readout, a
3-bit Euclidean code reproduces 0.96 of the exact top-1 selections at
$128\times$ compression, the bank fits in 6--\qty{33}{\mega\byte} per
layer, and its estimated in-array search energy is five to six orders
of magnitude below the energy of streaming it from HBM.
% ===== end 01 - Introduction.tex =====

%%%%%%%%%%%%%%%%%%%%%%%%%%%%%%%%%%%%%%%%%%%%%%%%%%%%%%%%%%%%%%%%%%%%%%
% ===== 02 - Background.tex =====
%%%%%%%%%%%%%%%%%%%%%%%%%%%%%%%%%%%%%%%%%%%%%%%%%%%%%%%%%%%%%%%%%%%%%%
\section{Background and Cost Model}
\label{sec:Background}

\subsection{The CAM design space}
\label{subsec:CAMspace}
\label{subsec:CAMtech}
Ordinary memory takes an address and returns content, so searching a
library means reading it row by row. A \cam compares every row
against the query at once, in place, and only reports matching
rows~\cite{pagiamtzis2006cam}: one access is one search of the whole
library, and the data that moves is the query rather than the bank.
A \textbf{binary \cam} (BCAM) cell holds one bit; a \textbf{ternary
\cam} (\tcam) cell holds a bit or a don't-care~\cite{ni2019fetcam}; a
\textbf{multi-bit \cam} (\mcam) cell holds one of several
levels~\cite{kazemi2021fefetcam}; an \textbf{analog \cam} (\acam)
cell holds a range~\cite{li2020acam}. A binary or ternary cell can
report only whether a position disagrees, so a row of them measures
Hamming distance. A multi-bit cell stores a level: in a multi-bit
\fefet cell the matchline current is approximately quadratic in the
difference between query and stored level, so $N$ cells on a
matchline sum to $\sum_i (Q_i-M_i)^2$, a squared Euclidean distance
with no multiplier
anywhere~\cite{kazemi2021fefetcam,kazemi2022hdcfefet}.\footnote{Squared
Euclidean stands in for Euclidean; the square root is monotonic.} The
array reports an \textbf{exact}
match, the \textbf{best} match (fewest mismatches), or a
\textbf{threshold} match (every row within a limit $\tau$); the
expected answer in monitoring is almost always empty, which threshold
mode reports without ranking anything. The cheapest cells (BCAM at
$\sim$6--10 transistors in any CMOS process; \tcam at $\sim$16T, or
two \fefets~\cite{pagiamtzis2006cam,ni2019fetcam}) are also the
Hamming-only ones, so an encoding that holds accuracy under a Hamming
metric can run on the cheapest cell available;
Sec.~\ref{subsec:BitWidth} compares that against a multi-bit array.

SRAM-based cells are buildable in any CMOS process today; \fefet-,
RRAM-, and MRAM-based cells trade that ubiquity for density and
nonvolatility, with device work moving toward the back end of
line~\cite{dutta2020m3dfefet}, where a static signature bank could
eventually sit above the compute that produces the activations. None
of these would be built for one application -- an SRAM associative
processor already performs Hamming search at billion-vector scale
commercially~\cite{gsi2025apu}, and \cam accelerators have been
proposed for one-shot learning~\cite{ni2019fetcam}, nearest-neighbor
search~\cite{kazemi2021fefetcam}, and range
matching~\cite{li2020acam} -- so the contribution here is the
workload mapping and the library-size analysis, not a new cell.

\subsection{Preliminaries and related work}
\label{subsec:Probes}
\label{subsec:RelatedWork}
\textbf{Probes:} A token is carried through a transformer by a
running vector, the \emph{residual stream}; after layer $\ell$ it is
$h^{(\ell)}$, $d$ numbers ($d=8{,}192$ for a 70B model). No single
coordinate means ``deception,'' and the simplest whole-vector
detector is a \emph{linear probe}, one stored vector $w$ fitted
offline on labeled activations and a threshold $\theta$, with
$w \cdot h > \theta \Rightarrow$
flag~\cite{marks2023geometry,burns2023ccs,zou2023repe}, costing
roughly $10^{-5}$ of a 70B forward pass.
\textbf{Grading:} \auroc is the fraction of (deceptive, honest) pairs
the detector orders correctly, with
\emph{cross-domain} evaluation -- fit on one domain, evaluate on
another~\cite{ying2026spectrum} -- as the stress test for
library size.\footnote{A \emph{domain} is a subject matter to be
dishonest about; a \emph{behavior} is a kind of dishonesty (lying
about a fact, flattering the user, underperforming on an
evaluation).}
\textbf{Related work:} The two-stage cascade
in~\cite{ccpp2026,anthropic2026ccblog} and the deployed Gemini
probes~\cite{gemini2026probes} are the closest published statements
of this workload; both operate at $K$ of order one. The
probe-fronted cascade reports 1--3.5\% overhead against 23.7\% for
the text-classifier pair it replaced~\cite{cc2025,ccpp2026}.
Approximate
nearest-neighbor search~\cite{charikar2002simhash,malkov2020hnsw}
visits a fraction of a bank, but a monitor cannot afford the misses
an adversary can aim at; an exhaustive scan has none to exploit.
\dt[still to verify: that the 1--3.5\% in ccpp2026 is measured over
the full cascade rather than the probe alone. Also: the two
``companion document'' pointers (red-team objections) are removed
from the body -- restore them if the appendix document goes to
reviewers]{}

%%%%%%%%%%%%%%%%%%%%%%%%%%%%%%%%%%%%%%%%%%%%%%%%%%%%%%%%%%%%%%%%%%%%%%
% ===== 03 - Library Scaling.tex =====
%%%%%%%%%%%%%%%%%%%%%%%%%%%%%%%%%%%%%%%%%%%%%%%%%%%%%%%%%%%%%%%%%%%%%%
\subsection{The cost model}
\label{sec:Scaling}
\label{subsec:CostModel}
\label{subsec:Ladder}
Let $K$ be the number of stored signatures per monitored layer, $d$
the residual width, $L_{\mathrm{mon}}$ the number of monitored
layers, $R$ the token rate, and $b$ the bytes per stored element. A
GPU reads the bank once per monitored token, while a \cam leaves the
bank in place and moves only the query, so its traffic is the
activation read that any monitor pays:
%
\begin{equation}
T_{\mathrm{GPU}} \;=\; K \, d \, b \, L_{\mathrm{mon}} \, R, \qquad
T_{\mathrm{CAM}} \;=\; d \, b \, L_{\mathrm{mon}} \, R
\quad\text{bytes/s,}
\label{eq:gputraffic}
\end{equation}
%
the latter independent of $K$. At $d=8{,}192$, fp16, one monitored
layer and $R=1{,}000$~tokens/s, monitoring traffic reaches an H100's
entire \qty{3.35}{\tera\byte\per\second}~\cite{nvidia_h100} at
$K\!\approx\!204$K, an A100's at $K\!\approx\!122$K, and a
Jetson-class part's
\qty{273}{\giga\byte\per\second}~\cite{nvidia_thor} at
$K\!\approx\!17$K. The crossings move with the token rate, a
deployment parameter rather than a property of the monitor, so
Fig.~\ref{fig:ladder} plots the limiting $K$ across rates from
batch-1 edge decoding (robot-control policies emit tens to $10^{3}$
action tokens/s with no batch to amortize
against~\cite{vlaperf2026}) to 70B-class serving at order $10^{3}$
tokens/s per datacenter GPU~\cite{lyceum2026tps}. Two properties of
Eq.~(\ref{eq:gputraffic}) matter. {\bf (1)} Latency is the wrong
way to judge a monitor's cost: at small $K$ the reads hide under
inference latency, so the monitor \emph{looks} free, but the bytes
and the energy still must move, per token. {\bf (2)} A per-token
cost multiplies with the token count -- the campaign runs above span
32 steps over roughly 20 hours and up to $10^{8}$
tokens~\cite{aisi2026cyber} -- so it determines whether an always-on
monitor stays on.
% ===== end 03 - Library Scaling.tex =====
% ===== end 02 - Background.tex =====

%%%%%%%%%%%%%%%%%%%%%%%%%%%%%%%%%%%%%%%%%%%%%%%%%%%%%%%%%%%%%%%%%%%%%%
% ===== 04 - Detector Mapping.tex =====
%%%%%%%%%%%%%%%%%%%%%%%%%%%%%%%%%%%%%%%%%%%%%%%%%%%%%%%%%%%%%%%%%%%%%%
\section{Mapping the Detector into a CAM}
\label{sec:Mapping}
Signatures can be mapped onto \cam arrays across the cell types of
Sec.~\ref{subsec:CAMspace}. The first mapping assumed the best-case
multi-bit device and is the reference point; Encodings 1 and 2 meet
the requirements below on hardware buildable today, while Encoding 3
is a bound rather than a \cam-resident design.

\subsection{What ``CAM-resident'' requires}
\label{subsec:Requirements}
For the monitor to be one memory structure rather than a memory plus
an accelerator: {\bf (1)} the query has to be producible from the
raw activation by something cheap, since whatever sits in front of
the array runs on every token; {\bf (2)} the comparison the array
performs natively has to be the comparison the detector needs, which
for a binary or ternary cell means a Hamming distance; and {\bf (3)}
the decision rule has to consume what a matchline can report -- a
winner, a set of rows under a threshold, or a count -- rather than a
full sorted score vector. Requirement~(1) has a methodological
consequence: a trained front end can learn to do the classification
itself, leaving the memory as decoration (Sec.~\ref{subsec:Audit}).

\subsection{The first mapping: learned projection, 3-bit Euclidean}
\label{subsec:FirstPipeline}
\label{subsec:PathBack}
A multi-bit, \fefet-class cell holds a level rather than a bit -- a
richer representation, or the same representation in fewer cells --
with the matchline computing the Euclidean distance the detector
actually uses; a down-projection is needed regardless
(requirement~{\bf (1)}). The first pipeline took the layer-33
activation, built a bank of mutually
distinct deception directions~\cite{ravfogel2020inlp}, compressed the
activation with a \emph{learned} projection from $d=8{,}192$ to
$d'=128$, quantized query and rows to 3 bits with learned-step
quantization~\cite{esser2020lsq}, and matched by Euclidean
winner-take-all. $d'=128$ matches the row
geometry of a multi-bit \fefet array -- 128 cells sharing a
matchline at 3 bits each -- so a stored row is
384-bit-\emph{equivalent}, the matched operating point of
Sec.~\ref{sec:Results}. The encodings below trade its
multi-bit cell for Hamming-distance solutions realizable by NVMs and
by plain CMOS; Sec.~\ref{subsec:BitWidth} evaluates the
recombination, a frozen projection on a Euclidean matchline.

\subsection{Encoding 1: sign random projection, ternary (or via other
NVMs)}
\label{subsec:SignRP}
Draw $P \in \mathbb{R}^{m \times d}$ once and freeze it, then encode
an activation as the sign pattern of its projection,
%
\begin{equation}
c_j(h) \;=\; \operatorname{sign}\!\big( (Ph)_j \big), \qquad j = 1 \ldots m,
\label{eq:signrp}
\end{equation}
%
with a dead zone $|(Ph)_j| < \epsilon$ mapped to don't-care, which
makes the code ternary and requires a wildcard-capable cell. This is
the standard sign-random-projection (SimHash)
construction~\cite{charikar2002simhash}: the Hamming distance between
two codes estimates the angle between the vectors, and it is an XOR
plus a population count, which a plain \cam array computes natively
(a code wider than an array row is split across subarrays and
aggregated, as in the measured \fefet - \cam
of~\cite{kazemi2022hdcfefet}).

\subsection{Encoding 2: thermometer coding (plain binary CMOS)}
\label{subsec:Thermometer}
Quantize each projected dimension into $s$ levels and encode each
level in unary, so level $q$ becomes $q$ ones followed by $s-q$
zeros. For two quantized vectors $u,v$,
%
\begin{equation}
d_{\mathrm{Ham}}\big(\mathrm{therm}(u), \mathrm{therm}(v)\big)
\;=\; \sum_i \big| u_i - v_i \big|,
\label{eq:thermo}
\end{equation}
%
i.e., Hamming distance on the code is exactly the $L_1$ distance on
the quantized values, so a cell that can only report disagreement
supports a magnitude-sensitive metric -- at roughly 6--10 transistors
rather than a \tcam's 16, since no wildcard is needed. A
1\fefet-1C array (a DRAM-like cell with the \fefet as access
transistor) has been extended to ternary \cam operation with a
fabricated proof of concept whose matchline voltage is linear in
Hamming distance~\cite{fefet1c2026}, so either encoding could also be
hosted there, nonvolatile.
\dt[DAC-thermosensitive.pdf: cite here with its plausible extension
once the PDF arrives from MN's work machine; also settle the
fefet1c2026 bib entry's authors before submission]{}

\subsection{The decision rule}
\label{subsec:Decision}
The stored rows are labeled example activations rather than fitted
directions: each domain stores several hundred to a few thousand
encoded activations, each labeled honest or deceptive. A query is
scored by the fraction
of its $k=31$ nearest stored rows, under the array's metric, that
carry the deceptive label (the gain flattens well before 31); \auroc
is computed on that score directly, and the rule consumes only what a
best-match or threshold matchline can report.

\subsection{Encoding 3, as a ceiling: shortlist plus exact rerank}
\label{subsec:TwoStage}
To bound what the compression costs, we also run a two-stage
variant: a thermometer-Hamming search of the whole bank keeps the
closest rows as a shortlist, and an exact cosine rerank on the raw
activations scores the shortlist. With the shortlist equal to the
whole bank, stage 2 alone defines the full-precision ceiling for the
same data and protocol; the construction is standard in large-scale
retrieval~\cite{faisscuvs2025}.
% ===== end 04 - Detector Mapping.tex =====

%%%%%%%%%%%%%%%%%%%%%%%%%%%%%%%%%%%%%%%%%%%%%%%%%%%%%%%%%%%%%%%%%%%%%%
% ===== 05 - Results.tex =====
%%%%%%%%%%%%%%%%%%%%%%%%%%%%%%%%%%%%%%%%%%%%%%%%%%%%%%%%%%%%%%%%%%%%%%
\section{Workload 1: Deception Detection}
\label{sec:Results}
We first evaluate the mappings on a deception detection workload,
where the library is set by how many behaviors and domains someone
decides to watch -- the $K\approx10^{2}$--$10^{3}$ regime of
Fig.~\ref{fig:ladder}.

\begin{figure*}[t]
\centering
\includegraphics[width=7.0in]{fig_grids_v3.pdf}
\caption{Cross-domain \auroc on the ten-domain grid,
Llama-3.3-70B-Instruct activations at layer 33; rows index the domain
stored in the bank, columns the domain of the held-out test
activations. (a) 384-bit \signrp ternary code (wildcard-capable
\tcam cell); (b) 384-bit thermometer code (plain binary \cam cell);
(c) Hamming shortlist over the whole bank followed by an exact cosine
rerank on the raw activations, i.e., the full-precision ceiling for
the same data and protocol.}
\label{fig:grids}
\end{figure*}

\subsection{Experimental setup}
\label{subsec:Setup}
We follow the cross-domain protocol of the Truthfulness Spectrum
study~\cite{ying2026spectrum}: a detector built from one domain's
labeled honest and deceptive activations is scored on every other
domain's. Activations are Llama-3.3-70B-Instruct at
layer 33~\cite{llama3herd}, the layer the source study tuned, so our
grids compare against theirs. The ten domains are definitional,
empirical, fictional, logical, ethical, sycophancy,
\texttt{repe\_honesty}~\cite{zou2023repe}, roleplaying, insider
trading, and sandbagging; within each domain 70\% of examples form
the stored bank and 30\% are held out, and banks range from 831 rows
(definitional) to 3{,}201 (sycophancy).

\subsection{The grids}
\label{subsec:Grids}
Fig.~\ref{fig:grids} shows the grids for the two single-stage
encodings and the exact-rerank ceiling (Table~\ref{tab:encodings}
summarizes them with the hardware each requires).

\textbf{In-domain fidelity is essentially unaffected by the
encoding:} Diagonal means are 0.920 (\signrp ternary, 384~b) and
0.919 (thermometer, 384~b) against 0.928 for the exact rerank: a
384-bit code in a plain array costs 0.008--0.009 \auroc against
full-precision cosine similarity over 8{,}192 values.

\textbf{Cross-domain fidelity is where compression shows:}
Off-diagonal means are 0.640 (ternary) and 0.579 (thermometer)
against 0.705, a cost of 0.065--0.126 \auroc. At matched bits ternary
leads by 0.061 on average and in 63 of 90 cells -- but a ternary
position may cost roughly two bits of hardware. At matched cell
count, the thermometer code at 768~b (256 directions at 3 bits)
scores 0.924 in-domain and 0.631 cross-domain: within 0.009 of
ternary's mean, ternary ahead in 54 of 90 cells, and the worst cell
improving from 0.259 to 0.319 -- at matched silicon the two appear
to be within noise on the mean, with ternary keeping a small
cell-wise edge.

\textbf{The pattern of failure reproduces the published one:} All
three grids keep the factual block high and put sycophancy near
chance from every other row, as the source study
reports~\cite{ying2026spectrum}; the full-precision grid has the
same holes in the same places, so the ceiling on transfer is set by
the representations rather than by our compression. Pooling does not
change this: a bank built from all ten
training splits (11{,}278 rows) averages 0.906 (thermometer), 0.909
(ternary), and 0.927 (ceiling) over the ten test domains, with
sycophancy still at 0.63, unlike the source study's combined probe at
0.97. Covering a behavior suite therefore means storing rows per
behavior and domain -- the library scale this paper prices.

\begin{table}[t]
\centering
\caption{The encodings, their measured fidelity on the ten-domain
grid, and what each needs from a cell. Bank size is for a
4{,}000-row exemplar bank at the given code width; $^{*}$matched
cell count against ternary at 384~b.}
\label{tab:encodings}
{\scriptsize\setlength{\tabcolsep}{3.5pt}
\begin{tabular}{@{}lccll@{}}
\toprule
\textbf{Encoding} & \textbf{in-dom.} & \textbf{cross-dom.} &
\textbf{Cell needed} & \textbf{Bank} \\
 & mean & mean & & (4K rows) \\
\midrule
\signrp ternary, 384\,b       & 0.920 & 0.640 & \tcam, $\sim$16T   & 192\,kB \\
Thermometer, 384\,b           & 0.919 & 0.579 & BCAM, $\sim$6--10T & 192\,kB \\
Thermometer, 768\,b$^{*}$          & 0.924 & 0.631 & BCAM, $\sim$6--10T & 384\,kB \\
Shortlist $+$ rerank          & 0.928 & 0.671 & BCAM $+$ FP unit   & 192\,kB $+$ raw \\
Exact rerank (ceiling)        & 0.928 & 0.705 & FP baseline        & 65\,MB \\
\bottomrule
\end{tabular}}
\end{table}

\subsection{Neither wider codes nor a richer cell closes the gap}
\label{subsec:BitWidth}
More code capacity does not close the gap, in either encoding: from
384 to roughly 2{,}048 bits, 23 of 90 off-diagonal
cells drop by more than 0.05 \auroc for ternary (worst $-0.303$,
sandbagging tested on insider trading) and 12 of 90 for thermometer
(worst $-0.209$), while the diagonal is flat or slightly up. A
richer cell does not help either once the projection is frozen: with
the same frozen projection, rounded to eight levels, a Euclidean
matchline (128 multi-bit cells per row)
scores 0.920 in-domain and 0.579 cross-domain against 0.920 and 0.593
for Hamming distance on the corresponding thermometer code (896
binary cells per row), ahead in only 37 of 90 cells. The first
pipeline's edge therefore came from its trained projection rather
than its matchline; with a frozen projection the multi-bit cell's
remaining benefit is fewer cells per row. On this grid, then, the
costlier cell does not buy accuracy, and transfer, if it comes, comes
from storing more vectors rather than wider ones -- the library
growth Sec.~\ref{subsec:CostModel} prices.

\subsection{A fixed shortlist does not travel across bank sizes}
\label{subsec:Shortlist}
The two-stage variant appears to recover the ceiling -- with an
absolute shortlist of 400 rows the diagonal mean is unchanged at
0.928 and the off-diagonal mean falls only to 0.671 -- but the
aggregate hides a failure of a \emph{fixed} parameter. Coverage, how
often the true full-precision top-1 neighbor lands inside the
Hamming shortlist, is 82.5\% for the 831-row definitional bank but
18.3\% for the 3{,}201-row sycophancy bank, so the sycophancy row
loses 0.204 \auroc on average and up to 0.468 in one cell.
Recovering that mean takes $k$ at 75--100\% of the bank, where it is
no longer a shortlist. A \cam-resident design should therefore use the
single-stage encodings, or a threshold rather than a top-$k$
criterion.

\subsection{The destroyed-vector audit}
\label{subsec:Audit}
\label{subsec:AuditResults}
A pipeline of the form ``trained front end, then memory'' can pass
its accuracy test while the memory contributes nothing, so we
re-measure every configuration with the stored rows replaced by
random vectors of the same shape; the accuracy that survives is the
\emph{front-end share}, where high is bad. Across five variants of
the first pipeline it was 98--100\% for a learned binary hash and
66--85\% for a weighted thermometer -- the two highest-scoring
variants, one of which reached 0.993 \auroc, both retired -- against
24--28\% for 3-bit LSQ Euclidean, 20\% for a direct sign code, and
4--14\% for a random sign projection, which is why Encodings~1 and~2
use frozen projections. Those encodings pass -- randomizing the bank collapses the 2{,}048-b
\signrp grid to a mean of 0.50 with 83\% of cells within 0.4--0.6,
and a label-shuffle control at 384~b gives 0.500 (ternary) and 0.501
(thermometer) -- and we report front-end share alongside \auroc for
every configuration.
\dt[IB: confirm the wording -- the 384-b numbers you sent are from the
label-shuffle audit, a different control from the randomized-bank
audit; the text states the two separately. If a randomized-bank run
exists at 384~b, swap it in]{}
% ===== end 05 - Results.tex =====

%%%%%%%%%%%%%%%%%%%%%%%%%%%%%%%%%%%%%%%%%%%%%%%%%%%%%%%%%%%%%%%%%%%%%%
% ===== 06 - Large K.tex =====
%%%%%%%%%%%%%%%%%%%%%%%%%%%%%%%%%%%%%%%%%%%%%%%%%%%%%%%%%%%%%%%%%%%%%%
\section{Workload 2: Whole-Vocabulary Readout}
\label{sec:LargeK}
\label{subsec:Jspace}
At the other end of the $K$ axis the library is set by the model's
vocabulary. The readout method of Sec.~\ref{sec:Introduction}
reports one internal pattern per vocabulary
token~\cite{anthropic2026workspace} (and ``manipulation'' active
while a model produced falsified data during a pre-release audit), so
reading this space means comparing the internal state against
roughly 128K stored signatures for a Llama-family tokenizer, at
several layers, every token; nobody chose that size, and phrases
would only enlarge it.

\subsection{An in-house reproduction, at watchlist scale}
\label{subsec:Watchlist}
We reproduced the readout in-house on Llama-3.1-8B; about 13 of 32
layers carry a usable readout, the source of the ``13--25 usable
layers'' factor above. Signatures are
precomputed offline, so a watched word costs one comparison per
token, and we stored 47 deception-adjacent signatures per monitored
layer as \cam rows. In a staged replacement scenario, five of ten
watch words -- ``scandal,'' ``evidence,''
``leaked,'' ``fake,'' ``shutdown'' -- reached rank~0 by exact
full-vocabulary rank at interpretable moments in the 2{,}048 prompt
tokens, ``threat'' rank~1 and ``blackmail'' rank~2, none present in
the prompt, while the reply itself was cooperative and unremarkable. These are preliminary numbers: one
model, one scenario, 47 signatures.\footnote{The fidelity values below were read from result
charts rather than result files; verification is queued.}

\subsection{Fidelity at \cam precision}
\label{subsec:JspaceFidelity}
The full-precision dot product against the stored patterns \emph{is}
the ground truth here, so no labels are needed: we record which of
the 47 stored vectors fire above a global threshold under the exact
computation on Wikitext activations, and ask how well each
\cam-compatible encoding reproduces that list under two readouts,
\emph{threshold} (every firing row reported; metric per-row firing
F1) and \emph{top-1} (the closest row reported; metric agreement with
the exact choice). Fig.~\ref{fig:jspace} plots both against storage
bits per row, the axis that sets the bank size; an fp16 row is
131{,}072 bits.

\begin{figure}[!t]
\centering
\includegraphics[width=3.4in]{fig_jspace_v2.pdf}
\caption{Whole-vocabulary readout reproduced through a \cam: (a)
per-row firing F1 and (b) top-1 agreement with the exact readout,
against storage bits per \cam row, for a bank of 47 vectors per
layer; the star marks an uncompressed fp16 row (131{,}072 bits).
Values digitized from the source result charts.}
\label{fig:jspace}
\end{figure}

\textbf{Two orders of magnitude of compression cost less than it
might:} A 3-bit Euclidean code at $\sim$$10^{3}$ bits per row, a
$128\times$ reduction, reaches 0.96 top-1 agreement and 0.81 firing
F1; a full-precision projection quantized only in width (the
exact-$L_2$ series) reaches 1.00 top-1 at 2{,}048 bits, so most of
the remaining gap is the quantizer rather than the projection.
\textbf{Top-1 is markedly cheaper than threshold readout:} At every
width the top-1 curve sits 0.15--0.2 above firing F1 -- a
matchline-mode result: best-match mode tolerates a cheaper cell than
threshold mode at equal fidelity.
\textbf{Hamming and thermometer codes fall short here:} Both plateau
near 0.74 top-1 and 0.49 firing F1 at $10^{3}$ bits, the opposite
ordering from Sec.~\ref{sec:Results}. The workloads ask different
questions -- separating two labeled classes there, reproducing a
ranked list here -- and since a frozen projection removed the
Euclidean advantage on the grid (Sec.~\ref{subsec:BitWidth}), the
readout's advantage appears to come from its \emph{trained}
projection as much as from the magnitude-preserving metric.

\subsection{What full readout would cost}
\label{subsec:ReadoutCost}
Priced at 70B geometry, an uncompressed bank of 128K signatures is
\qty{2.1}{\giga\byte} per layer; streaming it on every token at
1{,}000 tokens/s over 25 layers is roughly
\qty{53}{\tera\byte\per\second}. The useful answer stays small -- of order 25 patterns are
active at any one time -- so the operation is detection rather than
readout, the question a threshold-mode \cam answers natively. Three
things would make it tractable: {\bf (1)} a \emph{watchlist} of the
few hundred words known to be worth watching (the case demonstrated
above, three orders of magnitude smaller, but it catches only what
somebody thought to list); {\bf (2)} \emph{threshold matching}
instead of ranking; and {\bf (3)} \emph{compressing the stored side}
-- at 384--2{,}048 bits per row the whole 128K bank is
6--\qty{33}{\mega\byte} per layer and fits on a die, provided it does
not move. Only the first has been demonstrated here.
% ===== end 06 - Large K.tex =====

%%%%%%%%%%%%%%%%%%%%%%%%%%%%%%%%%%%%%%%%%%%%%%%%%%%%%%%%%%%%%%%%%%%%%%
% ===== 07 - Hardware Analysis.tex =====
%%%%%%%%%%%%%%%%%%%%%%%%%%%%%%%%%%%%%%%%%%%%%%%%%%%%%%%%%%%%%%%%%%%%%%
\section{Hardware Implications}
\label{sec:Hardware}
\label{subsec:Tax}
\label{subsec:CacheCliff}
\textbf{Monitoring as a share of the host's own traffic:} A 70B
model at fp8 reads roughly \qty{70}{\giga\byte} of weights per
token; $K=100$K rows at $d=8{,}192$ add \qty{1.6}{\giga\byte}, about
2\%. At the edge, an 8B model at 4-bit reads roughly
\qty{4}{\giga\byte} per token, and $K=25$K at $d=4{,}096$ adds
\qty{205}{\mega\byte}, about 5\%, rising to roughly 20\% by
$K=100$K -- the same point arrives one to two orders of magnitude
earlier in $K$ on parts with no HBM at any price and a
15--\qty{60}{\watt} envelope.
\textbf{The cache cliff:} A Jetson-class part has roughly
\qty{4}{\mega\byte} of shared L2~\cite{nvidia_thor}; 500 stored fp16
rows at $d=4{,}096$ -- the four campaign behaviors over 13 layers --
is \qty{4.1}{\mega\byte}, that L2 almost exactly; below the cliff
monitoring is close to free, above it the bank streams from LPDDR on
every token. At the measured code widths today's banks compress back
under the cliff (4{,}000 rows at 2{,}048 bits is about
\qty{1}{\mega\byte}), but at whole-vocabulary scale even the
compressed bank is out of cache.

\subsection{Where the search engine could sit, and on what}
\label{subsec:Placement}
\label{subsec:Devices}
A monitor co-resident on the inference GPU competes with live
inference for cache and bandwidth; a \cam or CXL-attached appliance matches in place but sends
the query over a board-level bus; a back-end-of-line search engine
above GPU compute would keep query travel to microns with a
nonvolatile static library, but no foundry PDK offers BEOL \fefet
today, so it is a 2030s option.
\dt[MN: confirm you want the BEOL clause in a DATE paper at all; it
is the most speculative sentence in the section]{}
The encodings decide the cell: a thermometer code needs only a
binary SRAM \cam in any CMOS process; the \signrp ternary code needs
wildcard capability, in CMOS \tcam at higher area or a two-transistor
\fefet cell~\cite{ni2019fetcam}; a Euclidean matchline needs a
multi-level device~\cite{kazemi2021fefetcam}, which the two workloads
value differently -- no accuracy gain on the deception grid once the
projection is frozen (Sec.~\ref{subsec:BitWidth}), a clear lead on
the readout workload (Sec.~\ref{subsec:JspaceFidelity}). Whether that
lead survives a common frozen projection decides whether an engine
serving both ends of Fig.~\ref{fig:ladder} needs two cell types or
one. Write endurance is not a constraint: the bank is rewritten per
model release, not per token, which is what makes an NVM bank
reasonable here.

\subsection{Energy and latency, against a measured GPU baseline}
\label{subsec:CostModelGap}
\label{subsec:NextMeasurements}
On the GPU, the dominant term at large $K$ is reading the bank out of
DRAM: at a measured 3.97~pJ/bit for HBM2~\cite{oconnor2017fgdram},
reading the \qty{2.1}{\giga\byte} whole-vocabulary bank once is
roughly 67~mJ per token per layer before any arithmetic. A measured
baseline sits above that: on an H100 PCIe, scanning the fp16 bank
with the $k=31$ top-$k$ included costs 1.22~ms and 277~mJ per token
at batch~1 and 21~$\mu$s and 4.7~mJ at batch~64 (at an achieved
1.6--1.8~TB/s, so the $K\!\approx\!204$K crossing of
Sec.~\ref{subsec:CostModel} is nearer 107K in practice). On the \cam
side the same search is $K$ rows of 2{,}048-bit code compared in
place, about $2.6\!\times\!10^{8}$ cell comparisons at the published
0.1--1~fJ per cell per search for CMOS and fabricated \fefet
arrays~\cite{pagiamtzis2006cam,ni2019fetcam,liu2022evacam}: roughly
0.03--0.3~$\mu$J per token per layer, five to six orders of magnitude
under the DRAM read. The gap is $64\times$ from the code width
(2{,}048 stored bits against 131{,}072) and
$\sim$$4\!\times\!10^{3}$ from the substrate.\footnote{The \cam figure is array-level (matchline and
sense-amp energy), excluding the $\sim$$10^{6}$-MAC query projection
and the mismatch-count reduction across subarrays; the 67~mJ
estimate excludes distance arithmetic and cache reuse, which the
measured rows include.} Compression also helps the GPU: a
population-count search (FAISS) over the
same 384-bit codes costs 28~$\mu$s and 2.1~mJ per token at batch~1
(0.45~mJ above idle) and 0.5~$\mu$s and 0.04~mJ at batch~64, flat in
$K$ from $10^{3}$ to 131K rows -- about 3\% of the part's per-token
budget at $10^{3}$~tokens/s. What the array offers over that
baseline is the search energy itself, 0.03--0.3~$\mu$J against
40~$\mu$J--2~mJ, and freedom from contention with live inference,
rather than bandwidth relief alone. Latency divides the same way:
\tcam arrays report 0.3--4.4~ns per
search~\cite{liu2022evacam,ni2019fetcam}, so a whole-bank search with
subarrays in parallel plausibly completes in tens of nanoseconds,
against 1.22~ms and 21~$\mu$s measured. Table~\ref{tab:costmodel}
collects these entries; array-level characterization with
Eva-CAM~\cite{liu2022evacam} at a stated node is what would replace
the estimates.
\dt[MN: framing check. As written, the popcount baseline makes the
CAM case rest on search energy and non-contention rather than on
bandwidth alone -- say if you want it framed differently. Numbers
from IB's H100 PCIe runs, Sep 2026]{}

\begin{table}[t]
\centering
\caption{The cost model at $K=1.3\!\times\!10^{5}$ rows, per token
per layer. \textsc{est}: first-order values from literature figures
(HBM at 3.97~pJ/bit~\cite{oconnor2017fgdram}; \cam cells at
0.1--1~fJ/search~\cite{pagiamtzis2006cam,ni2019fetcam,liu2022evacam});
\textsc{meas}: one H100 PCIe at batch 1 / batch 64.}
\label{tab:costmodel}
{\scriptsize\setlength{\tabcolsep}{3pt}
\begin{tabular}{@{}lccc@{}}
\toprule
\textbf{Metric} & \textbf{\cam, 2{,}048-b rows} & \textbf{GPU, fp16 rows} &
\textbf{GPU, 384-b codes} \\
\midrule
Energy / search & \makecell[c]{0.03--0.3\,$\mu$J\\\textsc{est}} &
  \makecell[c]{277 / 4.7\,mJ\\\textsc{meas} (67\,mJ \textsc{est})} &
  \makecell[c]{2.1 / 0.04\,mJ\\\textsc{meas}} \\
Latency / search & \makecell[c]{0.3--4.4\,ns per array\\\textsc{est}} &
  \makecell[c]{1.22\,ms / 21\,$\mu$s\\\textsc{meas}} &
  \makecell[c]{28 / 0.5\,$\mu$s\\\textsc{meas}} \\
Traffic / token & $d\,b\,L_{\mathrm{mon}}$ & $K\,d\,b\,L_{\mathrm{mon}}$ &
  $K\,(48\,\mathrm{B})\,L_{\mathrm{mon}}$ \\
Bank / layer & 33\,MB & 2.1\,GB & 6\,MB \\
\bottomrule
\end{tabular}}
\end{table}
% ===== end 07 - Hardware Analysis.tex =====

%%%%%%%%%%%%%%%%%%%%%%%%%%%%%%%%%%%%%%%%%%%%%%%%%%%%%%%%%%%%%%%%%%%%%%
% ===== 08 - Conclusions.tex =====
%%%%%%%%%%%%%%%%%%%%%%%%%%%%%%%%%%%%%%%%%%%%%%%%%%%%%%%%%%%%%%%%%%%%%%
\section{Conclusions}
\label{sec:Conclusions}
\label{subsec:Limitations}
\label{subsec:Conclusions}
\textbf{Limitations:} Nothing here competes with a deployed
classifier; the measured claims concern the fidelity and cost of
stored-signature matching, one component of the always-on stage of
Sec.~\ref{sec:Introduction}. Deception is defined as an output that
contradicts the known correct output; the grid results classify one
activation per response rather than per token, at one tuned layer of
one model. The hardware comparison has no contention term, and the
\cam entries of Table~\ref{tab:costmodel} are literature-based
estimates.
\dt[regenerate Fig.~\ref{fig:jspace} from result files before
submission -- the plotted values were digitized from the 07-30 deck]{}

Always-on behavioral monitoring of a language model is a
library-search problem whose coverage and cost are both set by $K$;
at $K=1$ the operation is nearly free, while whole-vocabulary readout
fixes $K\!\approx\!10^{5}$ per
layer by construction. Accuracy alone does not establish that a
memory-based detector is using its memory, and reporting front-end
share alongside \auroc is what moved this work onto frozen
projections. With them, a 384-bit code in a plain array holds
in-domain deception \auroc within 0.009 of full-precision cosine
similarity and gives up 0.065--0.126 cross-domain -- a gap that
neither wider codes, a multi-bit cell, nor a pooled bank closes, so
covering more behaviors means storing more rows -- and a 3-bit
Euclidean code reproduces 0.96 of a whole-vocabulary readout's top-1
selections at $128\times$ compression. That bank is
6--\qty{33}{\mega\byte} per layer and fits on a die, and comparing it
in place is estimated at four to six orders of magnitude below the
measured energy of streaming it from HBM, and two to nearly five
orders below a GPU population-count search over the same codes. What
remains is Eva-CAM characterization at a stated node and, upstream of
any hardware, establishing the minimum $K$ that adequate behavioral
coverage requires.
% ===== end 08 - Conclusions.tex =====

\bibliographystyle{IEEEtran}
\bibliography{refs}

% ---------- review-only appendix (With Comments build only) ----------
\reviewonly{%
\clearpage
\section*{\textcolor{aiviolet}{Length-cut log (not in the submission)}}
\begin{pnote}[aiviolet]
\textbf{Baseline (Sep 16):} v7 tree consolidated into this single
file; 16 pages clean (${\approx}$14.7 body $+$ 1.3 references) against
the DATE limit of 6 pages $+$ 1 for references. The full change history
through v7 lives in the archived driver one level up.

\textbf{Round 1 (Sep 16), structural, per IB's calls:} {\bf (1)}
Sec.~III folded into Sec.~II as the ``Cost model'' subsection (labels
kept, so every Sec.~III cross-reference still resolves). {\bf (2)}
Limitations folded into one paragraph at the head of the Conclusions.
{\bf (3)} Old VI-D (library-size sweep -- digitized data, three known
bugs) cut; Fig.~5(c) is now orphaned and flagged. {\bf (4)} Old VII-F
(the measurement list) cut; its cross-references rewritten, one
sentence of remaining work stays in the Conclusions. {\bf (5)} The
four-domain first-pipeline results paragraph in V-B cut (those numbers
were the undergraduates' pipeline, not the ten-domain runs). {\bf (6)}
Decision rule reworded per IB: the score is the fraction of deceptive
rows among the 31 nearest, and \auroc is computed on it directly.
{\bf (7)} All resolved review notes, the old commented-out IV-C, and
the per-section history comments removed from the source; the open
orange flags remain. No figure or table was changed in this round --
those are being decided one at a time.

\textbf{Round 2 (Sep 16), compression, per IB's ``keep going'':}
{\bf (1)} \emph{Floats:} Fig.~1 is panel (b) alone (fig\_ladder\_v3;
the $K$-vs-date panel and the anchor strip are gone, with the anchor
paragraph); Fig.~2 is the ten-domain row alone (fig\_grids\_v3, labels
renamed per IB); old Figs.~3, 4, 6 and Tables~IV, V cut, their numbers
kept in the text; the readout figure is two panels at single-column
width (fig\_jspace\_v2); the cost table is compacted to four rows and
gains a measured GPU column. Figure sources: date/gen\_cut\_figs.py.
{\bf (2)} \emph{Text:} every section rewritten to roughly 40\% of its
length in the draft's voice rules; equations, the encoding
definitions, the three-observation structure of Secs.~IV and V, and
every quoted number are kept. Sec.~II-B (technologies) is folded into
II-A; the bit-width subsection and the recombination result share one
subsection (IV-C); Sec.~VI is three blocks (share-of-traffic and
cache cliff, placement and cells, energy/latency).
{\bf (3)} \emph{IB's results folded in:} matched cell count (768-b
thermometer) in IV-B and Table~II; pooled bank in IV-B; 384-b
label-shuffle control in IV-E (flagged: it is a different control from
the randomized-bank audit); recombination (frozen RP, Euclidean vs
Hamming) in IV-C, echoed in V-C and VI-A; H100 PCIe baselines (fp16
bank and FAISS popcount) in VI-B and Table~VI, with a framing flag for
MN; decision rule wording; domain renames.
{\bf (4)} \emph{Dropped with their text:} schmidt2026interp,
gao2024sae, lieberum2024gemmascope, kim2023alscn (unverified author
list), apple2025foundation, indyk1998lsh, deepseekv3,
anthropic2025cheapmonitors, the QAT footnote, the L$\cdot$d aside, the
``not symmetric'' paragraph (closed by the matched-cell result), and
MN's II-B placeholder flag (the subsection it marked no longer
exists).
\end{pnote}
}

\end{document}
